\documentclass[11pt]{article}
\usepackage[margin=.9in]{geometry}
\usepackage{amsmath,amssymb,amsthm}
\usepackage[T1]{fontenc}
\usepackage{lmodern,microtype,hyperref}
\newtheorem{theorem}{Theorem}
\newtheorem{lemma}[theorem]{Lemma}
\newtheorem{proposition}[theorem]{Proposition}
\newtheorem{corollary}[theorem]{Corollary}
\theoremstyle{remark}\newtheorem{remark}[theorem]{Remark}
\newcommand{\norm}[1]{\left\lVert#1\right\rVert}
\DeclareMathOperator{\diag}{diag}\DeclareMathOperator{\ran}{ran}
\newcommand{\R}{\mathbb R}\newcommand{\C}{\mathbb C}
\begin{document}
% No preamble. Notebook 17 joint-tail module. Unique labels: jt-.
% Requires amsmath, amssymb, amsthm and a one-argument \norm macro.

\section{Joint Gaussian bounds for prescribed bad blocks}
\label{jt-section}

The input to this section is a counting-coordinate covariance
operator bound. A bound on each coordinate variance alone does
not suffice. We first give a real Gaussian theorem, including
degenerate covariances and overlapping tested sets. We then prove
a normalized complex Gaussian extension under an explicit
relative quadratic hypothesis. All statements are initially
finite-dimensional; their constants do not depend on the ambient
dimension or the number of tested blocks.

\subsection{The covariance input and microscopic blocks}
\label{jt-input}

Let $Y=(Y_i)_{i\in I}$ be a centered real Gaussian vector with
counting-coordinate covariance
\begin{equation}\label{jt-covariance}
 0\le\Sigma\le K I,\qquad K>0.
\end{equation}
For each block $b$ let $I_b\subset I$ be a nonempty tested set.
Assume fixed constants $m,\nu$ satisfy
\begin{equation}\label{jt-incidence}
 |I_b|\le m,\qquad
 \sup_i\#\{b:i\in I_b\}\le\nu.
\end{equation}
Thus $m$ is the number of scalar tests per block and $\nu$ is
their overlap multiplicity. Define
\begin{equation}\label{jt-bad-event}
 \mathsf B_b(R)=\{\max_{i\in I_b}|Y_i|>R\},\qquad R>0.
\end{equation}
Tests may be strongly correlated. The overlap counts repeated
$Y$ coordinates, not overlap of the underlying field variables
from which a derivative test was formed. Enlarging a block to
include a fixed stencil changes $m$ and $\nu$ by fixed factors.

\begin{theorem}[Joint real Gaussian bad-block bound]
\label{jt-real-theorem}
For every finite set $S$ of $k$ blocks,
\begin{equation}\label{jt-real-bound}
 \mathbb P\!\left(\bigcap_{b\in S}\mathsf B_b(R)\right)
 \le\left[
 m\sqrt2\exp\!\left(-\frac{R^2}{4\nu K}\right)
 \right]^k.
\end{equation}
The prefactor $m\sqrt2$ may instead be $2^{m/2}$, and the
right side may be replaced by its minimum with one.
\end{theorem}

\begin{proof}
Choose one coordinate $i_b\in I_b$ for each block. The vector
$W=(Y_{i_b})_{b\in S}$ has covariance at most $\nu K I$:
the coordinate-selection map, with repetitions allowed, has
squared operator norm at most $\nu$. Diagonalizing this
nonnegative covariance, allowing zero eigenvalues, gives
\begin{equation}\label{jt-square-mgf}
 \mathbb E e^{t|W|^2}
 =\det(I-2t\Sigma_W)^{-1/2}
 \le(1-2t\nu K)^{-k/2},\qquad 0\le t<(2\nu K)^{-1}.
\end{equation}
If every chosen coordinate exceeds $R$ in absolute value,
$|W|^2>kR^2$. Take $t=(4\nu K)^{-1}$ in Markov's inequality.
The resulting bound is $2^{k/2}e^{-kR^2/(4\nu K)}$.
There are at most $m^k$ coordinate choices, and their union
contains the event in \eqref{jt-real-bound}.

Alternatively collect every tested coordinate in the selected
blocks, retaining repetitions. This vector has at most $mk$
entries and covariance at most $\nu K I$. On the all-bad event
its squared norm is at least $kR^2$. Applying
\eqref{jt-square-mgf} with dimension $mk$ proves the alternative
prefactor. Neither proof assumes independent blocks.
\end{proof}

If a local norm obeys
$N_b(Y)\le c_b\max_{i\in I_b}|Y_i|$, its bad event is
bounded by the theorem with threshold $R/c_b$. For a Euclidean
block norm, the second proof applies directly to
$\sum_{i\in I_b}Y_i^2>R^2$, with prefactor $2^{m/2}$.
For $gY$ replace $R$ by $R/g$, equivalently $K$ by $g^2K$.

\subsection{Exponential squares, entropy, and connected sets}
\label{jt-entropy-section}

Put $E_b(Y)=\sum_{i\in I_b}Y_i^2$. For nonnegative $t_b$ of
finite support define
\[
 a_i=\sum_{b:i\in I_b}t_b,\qquad A=\operatorname{diag}(a_i).
\]
Then $\norm A\le\nu t_{\max}$ and
$\operatorname{Tr}A\le m\sum_b t_b$. The exact Gaussian
determinant, applied to $\Sigma^{1/2}A\Sigma^{1/2}$, gives
\begin{equation}\label{jt-exponential-square}
 \boxed{\quad
 \log\mathbb E\exp\!\left(\sum_b t_bE_b(Y)\right)
 \le\frac{Km\sum_b t_b}{1-2K\nu t_{\max}},
 \qquad 2K\nu t_{\max}<1.
 \quad}
\end{equation}
Indeed each eigenvalue is at most $K\nu t_{\max}$, their
sum is at most $K\operatorname{Tr}A$, and
$-\log(1-x)\le x/(1-x_{\max})$ for $0\le x\le x_{\max}<1$.
This proof also permits singular covariance.

For a probability law $\rho$ absolutely continuous with respect
to the Gaussian law $\mu$, the entropy variational inequality
therefore gives
\begin{equation}\label{jt-relative-entropy}
 \sum_b t_b\mathbb E_\rho E_b(Y)
 \le D(\rho\Vert\mu)
       +\frac{Km\sum_b t_b}{1-2K\nu t_{\max}}.
\end{equation}
Truncation gives the same assertion with extended expectations
when necessary.

Let $p(R)$ be the bracket in \eqref{jt-real-bound}, decreased
to at most one if desired. Expanding a finite product with
$z_b\ge0$ yields
\begin{equation}\label{jt-bad-count-product}
 \mathbb E\prod_b(1+z_b\mathbf1_{\mathsf B_b(R)})
 \le\prod_b(1+p(R)z_b).
\end{equation}
In particular, for $N$ tested blocks and $s\ge0$, the
exponential moment of their bad-block count is at most
$(1+p(R)(e^s-1))^N$. This is a moment statement, not an
assertion that the process is independent or stochastically
dominated for every increasing event.

Give the blocks an adjacency graph of maximum degree at most
$\Delta\ge1$. The number of connected $k$-block sets containing
a specified root is at most $(4\Delta)^{k-1}$. One way to see
this is to choose a canonical rooted traversal tree: ordered
rooted tree shapes cost at most $4^{k-1}$, and each edge has
at most $\Delta$ neighbor choices. Overcounting only enlarges
the bound. Consequently, whenever $4\Delta e^a p(R)<1$,
\begin{equation}\label{jt-connected-sum}
 \boxed{\quad
 \mathbb E\sum_{\substack{S\ni b_0\\S\ {
m connected}}}
 e^{a|S|}\prod_{b\in S}\mathbf1_{\mathsf B_b(R)}
 \le\frac{e^a p(R)}{1-4\Delta e^a p(R)}.
 \quad}
\end{equation}
This is a summable Gaussian bad-set estimate. It does not
estimate an interacting polymer activity or prove its strict
background dependence.

\subsection{The homogeneous flat applications}
\label{jt-flat-applications}

For the previously constructed d4 constrained flat Gaussian,
the covariance under the physical pairing is $\mathcal C$,
whereas its coordinate covariance is
$\operatorname{Cov}_{\rm count}(Z)=\eta^{-4}\mathcal C$.
The dimensionless derivative vector therefore has exactly
\begin{equation}\label{jt-derivative-covariance}
 \operatorname{Cov}_{\rm count}(\eta^2\nabla Z)
 =\nabla\mathcal C\nabla^T\le K_D I.
\end{equation}
The inequality follows from the proved real order
$\mathcal C\le\Gamma\le G_0$ and the uniform first-derivative
sandwich of $G_0$. Input and output fine spaces have the same
physical scalar measure, so the operator norm of this coordinate
matrix is unchanged on passing to counting Hilbert norms.
The color basis is real trace-orthonormal.

Fixed microscopic blocks contain a fixed number of tests, so
Theorem~\ref{jt-real-theorem} and
\eqref{jt-exponential-square}--\eqref{jt-connected-sum} apply
uniformly over the full homogeneous refinement. Rescaled curls
are fixed local combinations of this derivative vector and
have the same conclusion with fixed norm and incidence factors.
The sampled microscopic H\"older term satisfies
$\eta^{2+\alpha}[\nabla Z]_\alpha
\le C\norm{\eta^2\nabla Z}_\infty$ and can be included at
fixed cost. A unit physical cell contains $M^4$ fine sites,
so treating it as one block here would make $m$ grow with $M$.

For the full macroscopic Gaussian the potential vector instead
has covariance
\begin{equation}\label{jt-potential-covariance}
 \operatorname{Cov}_{\rm count}(\eta Z)=\eta^{-2}\mathcal C.
\end{equation}
Its counting-operator norm is not uniformly bounded over
refinement. At zero coarse momentum a transverse fine alias
$r=e_1$, polarized in direction $2$, lies in the constraints
and has Maxwell eigenvalue $4M^2\sin^2(\pi/M)\le4\pi^2$.
Thus \eqref{jt-potential-covariance} has norm at least
$M^2/(4\pi^2)$. Pointwise potential variance bounds do not
remove this operator obstruction.

For the first bare homogeneous step of a fixed ratio $L$, the
exact physical scaling established in the preceding section
instead gives the joint value/derivative covariance
\begin{equation}\label{jt-fixed-step-covariance}
 \Sigma_{\rm step}=S_L\mathcal C_LS_L^T\le K_LI,
 \qquad S_L=\begin{pmatrix}LI\\\nabla_{1/L}\end{pmatrix}.
\end{equation}
Here $L$ is fixed and the number of retained cells is arbitrary.
The same theorem therefore applies to both components together,
including their cross covariances. The real reduced-Hessian
comparison of that bare step enlarges $K_L$ by at most
$(1-\epsilon)^{-1}$ for chart-coordinate tests. At real
backgrounds it also gives a uniform bound for the separately
identified linearized-potential tests, so the real theorem
applies to those as well.

These applications do not replace the precision in a later
effective step by the bare fixed-$L$ Maxwell form. A later
coarse precision and its separate current terms require their
own covariance identification and uniform bound.

\subsection{A normalized complex quadratic extension}
\label{jt-complex-section}

Whiten the real baseline Gaussian on its nondegenerate slice:
let $\zeta$ be standard real Gaussian in $\mathbb R^N$ and
let $Y=A\zeta$, with $AA^T\le KI$. The observable covariance
may be singular. Suppose $E$ is complex symmetric and
\begin{equation}\label{jt-relative-perturbation}
 E^T=E,\qquad \norm E_{2\to2}\le\epsilon\le\frac14.
\end{equation}
Define the normalized complex Gaussian functional
\begin{equation}\label{jt-complex-measure}
 d\mu_E(\zeta)=\det(I+E)^{1/2}
      e^{-\zeta^TE\zeta/2}\,d\gamma_N(\zeta),
\end{equation}
with the analytic determinant branch equal to one at $E=0$.
Its total mass is one. For a physical quadratic perturbation
$V_X$ of the positive baseline precision $K_X$, condition
\eqref{jt-relative-perturbation} means exactly
$\norm{K_X^{-1/2}V_XK_X^{-1/2}}\le\epsilon$.
An absolute smallness assertion about matrix entries is not a
substitute for this relative bound. At complex $E$ the functional
is not a positive probability law.

\begin{theorem}[Complex Gaussian local-event bound]
\label{jt-complex-theorem}
Under \eqref{jt-incidence} and
\eqref{jt-relative-perturbation}, every set $S$ of $k$ blocks
satisfies
\begin{equation}\label{jt-complex-bound}
 \boxed{\quad
 \left|\mu_E\!\left(\bigcap_{b\in S}\mathsf B_b(R)\right)\right|
 \le\left[2^m\exp\!\left(-\frac{R^2}{6\nu K}\right)\right]^k.
 \quad}
\end{equation}
There is no ambient-$N$ factor. The events use fixed real linear
tests $Y=A\zeta$ and fixed real thresholds.
\end{theorem}

\begin{proof}
All tested coordinates in $S$ span a real row space of dimension
$r\le mk$. Make a real orthogonal change of the whitened
coordinates so that these tests depend only on $x\in\mathbb R^r$;
denote the untested orthogonal coordinates by $z$. Integrating
$z$ exactly gives a normalized complex Gaussian marginal on
$x$ with precision
\begin{equation}\label{jt-marginal-precision}
 B=I+E_{11}-E_{12}(I+E_{22})^{-1}E_{21}.
\end{equation}
Therefore
\begin{equation}\label{jt-marginal-control}
 \norm{B-I}\le\epsilon+\frac{\epsilon^2}{1-\epsilon}
       =\frac{\epsilon}{1-\epsilon}=:s\le\frac13,
 \qquad \operatorname{Re}B\ge(1-s)I=:\beta I,
 \qquad \beta\ge\frac23.
\end{equation}
The real-part assertion is an inequality for the real quadratic
form, equivalently the Hermitian numerical range of the complex
symmetric matrix. Since
$|\det B|\le\norm B^r\le(1+s)^r$, the modulus of the
normalized marginal density is bounded by
\begin{equation}\label{jt-marginal-domination}
 \left(\frac{1+s}{\beta}\right)^{r/2}
                d\gamma_{\beta^{-1}I_r}(x)
 \le 2^{r/2}d\gamma_{\beta^{-1}I_r}(x).
\end{equation}

Under this dominating real Gaussian the tested covariance norm
is at most $K/\beta$. Collect all tests in the selected blocks,
including repetitions. Its dimension is at most $mk$, overlap
at most $\nu$, and covariance norm at most $\nu K/\beta$.
On the all-bad event its squared norm exceeds $kR^2$.
Equation~\eqref{jt-square-mgf}, with
$t=\beta/(4\nu K)$, bounds its probability by
$2^{mk/2}\exp[-\beta kR^2/(4\nu K)]$.
Multiply by $2^{r/2}\le2^{mk/2}$ and use $\beta\ge2/3$.
This proves \eqref{jt-complex-bound}, including degenerate
observable covariance and the case $r=0$.
\end{proof}

The proof also bounds a bounded measurable factor depending
only on the same tested marginal, multiplied by the all-bad
indicator, by its supremum times the right side of
\eqref{jt-complex-bound}. It does not bound the integral of
the modulus of the full $N$-dimensional density over that
event without an $N$-dependent factor. Untested variables must
be integrated out before taking the modulus. An observable
using further coordinates requires enlarging the marginal
and accounting for its dimension.

For fixed real-coordinate events, the normalized integral is
holomorphic in $E$ on the complex operator ball. At each finite
$N$, the real-part lower bound gives Gaussian domination on
compact subballs and the determinant branch is analytic. If
$E$ depends holomorphically on sources in a known neighborhood
remaining inside \eqref{jt-relative-perturbation}, Cauchy on
that neighborhood gives factorial source derivatives of these
local event integrals, retaining the displayed bad-block decay
on a slightly larger parameter domain. This requires a proved
relative perturbation domain; it does not produce one for an
unmatched source Hessian.

In particular an indicator around tests that themselves move
with a complex background is not covered by this holomorphy
argument. For the reduced first-step Hessian, the fixed real
chart-coordinate tests are admissible. Background-dependent
linearized-potential tests are covered by the real covariance
theorem at real backgrounds, but require a separate analytic
cutoff construction at complex backgrounds.

Expansion of a finite product also gives
\begin{equation}\label{jt-complex-product}
 \left|\mu_E\prod_b(1+z_b\mathbf1_{\mathsf B_b(R)})\right|
 \le\prod_b(1+p_{\rm c}(R)z_b),\qquad
 p_{\rm c}(R)=2^m e^{-R^2/(6\nu K)},\qquad z_b\ge0.
\end{equation}
One must not cap $p_{\rm c}$ at one merely because the total
mass is normalized: complex expectations of events are not
probabilities. The same connected-set counting as
\eqref{jt-connected-sum} bounds the sum of absolute local
event expectations whenever its entropy inequality holds.

\subsection{Source match and application boundary}
\label{jt-scope}

For the homogeneous flat derivative application, the action,
coupling and physical pairing are the deposited conventions at
Raw I:464--479, I:12756--12760 and I:15643--15648. The flat
corrected average and its homogeneous composition are matched
at I:24349--24373 and I:47046--47048; the scalar projection is
matched at I:23937--23944. The analytic input actually used is
the established variational $\mathcal C$ and full-$G_0$
first-derivative sandwich. These are checks against the pinned
deposited transcription, not a new original-paper verification.

The real theorem therefore applies to that explicit homogeneous
flat constrained Gaussian, and to the fixed-ratio first bare
step with the covariance identification already proved. The
complex theorem applies under its explicit relative quadratic
hypothesis; for the reduced first bare step this hypothesis is
supplied by the preceding Hessian calculation after shrinking
the background ball. Later effective current-dependent steps
require their own identification.

These joint bounds are stronger than marginal variance or a
single union bound. They still do not integrate nonlinear
interactions, Haar or chart Jacobians, or regions outside the
chart. They do not supply a complete retained-history covariance,
an interacting polymer expansion, or a quantum continuum limit.

\end{document}
